% ERL Letter: statements (D-107). Data and code availability follow the IOP research data
% policy; the AI-use statement follows IOP's policy on generative AI (docs/journal_specs.md) and
% is a draft for the PI. Acknowledgements, funding and conflicts of interest are for the PI.
% Not counted in the main-text word count (tools/word_count_erl.py reports them separately).

\section*{Data availability statement}
The data that support the findings of this study will be openly available in a public repository with a DOI on publication\todo{PI: name the repository (e.g.\ Zenodo) and add its DOI; the release follows docs/DATA\_RELEASE.md}: the level-binned observed extents and waterlines of the North Caspian (1992--2026) and of the two arms of Great Salt Lake, the hindcast tables, the waterline elevation model with its documented accuracy, and every derived table behind the figures and numbers of this Letter. The input data (Landsat, the altimetry levels, the five bathymetry grids, ICESat-2, GEDI and Copernicus DEM heights, ERA5, OpenStreetMap under the Open Database License, ESA WorldCover, and the USGS gauge levels and lidar for Great Salt Lake) are openly available from the providers cited in the Methods and the Supplementary Information. Gauge levels of the Caspian Coordinating Committee on Hydrometeorology (CASPCOM) were used to estimate the offset to the Baltic Height System; the catalogue states no licence, so only the offset, its uncertainty and the code that recomputes it from the catalogue are released\citep{caspcom2026}. Results derived from the Copernicus DEM were produced using Copernicus WorldDEM-30 \textcopyright{} DLR e.V. 2010--2014 and \textcopyright{} Airbus Defence and Space GmbH 2014--2018 provided under COPERNICUS by the European Union and ESA; all rights reserved.

\section*{Code availability}
All code that produces the results, figures, tables and every number in this Letter, from the raw inputs to the text, is in the project repository under the MIT licence; one command (\texttt{make all}) rebuilds it. The repository, with its pre-registration, decision log and amendments, will be deposited in a public archive with a DOI; the PI decides when\todo{PI: repository URL and archive DOI at submission}.

\section*{Use of generative AI}
\todo{PI: review this draft statement against IOP's current policy on generative AI before submission.} This study made extensive use of generative AI. Agents based on large language models (Claude, Anthropic) worked under the direction of the authors: they wrote and ran the analysis code, the figure scripts and the tests, searched for and screened the literature, and drafted the pre-registration, the decision log and the text of this Letter and its Supplementary Information. The authors set the questions, approved the pre-registration and every decision recorded in the decision log, chose the venue, and reviewed the outputs. Every number in the text is generated by script from the results file, and references are retrieved, read and checked against their DOI records, with an audit trail kept in the repository. The AI tools are not authors, and the authors take full responsibility for the content.

\section*{Acknowledgements}
\todo{PI: acknowledgements.}

\section*{Funding}
\todo{PI: all funding sources with grant numbers, or a statement that the work received no specific funding.}

\section*{Conflict of interest}
\todo{PI: conflict-of-interest statement.}
